\subsection{包含关系}

定义1. 用以下记号表示的关系 \((\forall z)((z \in x) \Longrightarrow (z \in y))\) ，其中只出现字母 \(x\) 和 \(y\) ，按以下方式之一书写： \(x \subset y\) , \(y \supset x\) , `` \(x\) 包含于 \(y\) '', `` \(y\) 包含 \(x\) '', `` \(x\) 是 \(y\) 的子集''。关系“不 \((x \subset y)\) '' 写作 \(x \not\subset y\) 或 \(y \not\supset x\) .

按照第一章第1节第1号中提到的约定，这个定义蕴含以下元数学约定。设T和U为汇编；若我们在汇编 \(x \subset y\) 中将T替换为x、将U替换为y，我们得到一个汇编，记为 \(T \subset U\) ；如果我们用 x，y 表示与 x，y 不同且彼此不同的字母，并且它们既不出现在 T 中也不出现在 U 中，则汇编 \(T \subset U\) 恒同于 \((T|x)(U|y)(x|x)(y|y)(x \subset y,)\) ，因此由 CS8、CS9（第一章，§4，第 1 号）以及 CS5（第一章，§1，第 2 号），它恒同于 \((\forall z)((z \in T) \Longrightarrow (z \in U))\) ，前提是 z 是一个既不出现在 T 中也不出现在 U 中的字母。

从现在起，每当我们陈述一个数学定义时，我们将不再提及它所蕴含的元数学约定。

CS12。设 \(\pmb{T},\pmb{U},\pmb{V}\) 为组装，并设 \(\pmb{x}\) 是一个字母。则汇编 \((\pmb{V}|\pmb{x})(\pmb{T}\subset \pmb{U})\) 恒同于 \((\pmb{V}|\pmb{x})\pmb{T}\subset (\pmb{V}|\pmb{x})\pmb{U}\) .

这直接从 CS9（第一章，§4，第 1 号）和 CS5（第一章，§1，第 2 号）得出。

CF13。如果 \(\pmb{T}\) 并且 \(\pmb{U}\) 是项， \(\pmb{T} \subset \pmb{U}\) 是一个关系。

这直接从CF8（第一章，§1，第4节）得出。

¶ 每个形如 \(T \subset U\) （其中 T 和 U 是项）的关系称为包含关系。

从现在起，我们不再明确写出替换准则以及应在定义之后给出的构成准则。但应注意，这些准则在证明中常常会被隐式使用。

为了证明该关系 \(x \subset y\) 在一个理论中 \(\mathfrak{G}\) ，根据C27（第一章，§4，第1号），只需证明 \(z \in y\) 在通过附加 \(z \in x\) 对于公理 \(\mathfrak{G}\) ，其中 \(z\) 是一个与 \(x, y\) 以及理论的常数。在实践中我们说“令 \(z\) 为 \(x\) ''，并且我们试图证明 \(z \in y\) .

命题1。 \(x \subset x\) .

显然。

命题2。 \((x\subset y\) 并且 \(y\subset z)\Longrightarrow (x\subset z).\)

添加假设 \(x \subset y, y \subset z\) ，以及 \(u \in x\) 。则关系

\[
(u \in x) \Longrightarrow (u \in y), \quad (u \in y) \Longrightarrow (u \in z),
\]

为真，因此该关系 \(u \in z\) 为真。

